\documentclass[10pt,a4paper]{amsart}
\usepackage[margin=19mm]{geometry}
\usepackage{amsmath,amssymb,mathtools}
\usepackage[scr=rsfs,cal=euler]{mathalfa}
\usepackage{xfrac}
\usepackage[unicode,hidelinks,bookmarks=false]{hyperref}
\newcommand{\ie}{i.e.\ }
\newcommand{\cf}{cf.\ }
\newcommand{\resp}{respectively\ }
\newcommand{\Subsection}{\subsection}
\providecommand{\nobreakdash}{\nobreak\hbox{-}}
\setlength{\emergencystretch}{2em}
\multlinegap0pt
\numberwithin{equation}{section}
\newtheorem{theorem}{Theorem}[section]
\newtheorem{lemma}[theorem]{Lemma}
\newtheorem{corollary}[theorem]{Corollary}
\newtheorem{definition}[theorem]{Definition}
\newtheorem{remark}[theorem]{Remark}
\allowdisplaybreaks
\title[The constrained Lojasiewicz-Simon inequality on the unit sphere of a]{The constrained Lojasiewicz-Simon inequality on the unit sphere of a Hilbert space}
\author{Zdzis\l{a}w Brze\'zniak, Javed Hussain}
\date{Source: Electronic Journal of Differential Equations, Vol. 2026 (2026), No. 13, pp. 1--13; DOI 10.58997/ejde.2026.13; published February 11, 2026; publisher version}
\begin{document}
\maketitle
\noindent\emph{Electronic Journal of Differential Equations}, Vol. 2026 (2026), No. 13, pp. 1--13.\newline
ISSN: 1072-6691. DOI \href{https://doi.org/10.58997/ejde.2026.13}{10.58997/ejde.2026.13}. This work is licensed under a CC BY 4.0 license.\par\smallskip
\noindent\textit{Editorial scope.} Counted claim: original Theorem 3.3 of the published article (source label Theorem, native label Lwoj, source lines 543--553, printed as Theorem 3.3 with the parenthetical title Lojasiewicz-Simon inequality) together with its complete original author proof at source lines 555--701, delivered as the single primary proof body. No proof marker is added and no mathematics is written, repaired or re-derived. Necessary complete context is retained uncounted: the entire native text from the first section heading at source line 60 through the introductory sentences immediately preceding the counted statement at source line 542, that is the whole Introduction, the whole of Section 2 on functional settings, assumptions and preliminaries, and the whole first part of Section 3 up to the counted statement (Theorem 3.1 on precompactness with its proof and Corollary 3.2 with its proof), together with the subsection heading and the two introductory sentences of the Lojasiewicz-Simon subsection. Because the retained context is exactly the published text in published order and nothing inside the delivered ranges is omitted, every printed section, theorem, corollary, lemma, definition, remark and equation number of the delivered unit coincides with the published PDF. The delivered unit stops after the counted proof; the later Sections 4, 5 and 6 of the article are not delivered, and the citations of the delivered text are re-typeset from the published bibliography entries with the numbering of the published PDF. Printed slips kept literal: the source prints the Lojasiewicz name through an accent macro and the PDF prints the same accented form; the delivered text reproduces those source occurrences unchanged, including the one in the subsection heading at source line 536, and no printed word, symbol or numbering of the counted statement or of its proof was altered. No figure, image, table or external asset occurs anywhere in the delivered ranges, so nothing had to be omitted and no bounded source comparison was needed.
\section{Introduction}\label{sec-Introduction}

In this paper, we focus on the long-time dynamics of the projected nonlinear heat
equation
\begin{equation}\label{1.6}
\begin{gathered}
\frac{\partial u}{\partial t} = \pi_{u} \left( \Delta u - |u|^{2n-2} u \right), \\
u(0) = u_{0},
\end{gathered}
\end{equation}
posed in a bounded domain $\mathcal{O}\subset\mathbb{R}^{d}$ with homogeneous Dirichlet
boundary conditions. Our primary objectives are to establish the convergence of
solutions to equilibria, derive precise decay estimates, and prove the existence of
a compact global attractor.

The nonlinear term $|u|^{2n-2}u$ arises naturally as the variational derivative of
a polynomial potential and appears in a range of diffusion–reaction models.
The unit-sphere constraint in $\mathcal L^{2}(\mathcal O)$ enforces preservation of
a global invariant, such as mass, probability normalization, or total intensity,
depending on the modeling context.
Flows of this type occur, for instance, in normalized diffusion processes, constrained
relaxation dynamics, and evolution problems where the state is required to remain on
a prescribed manifold.
We are mainly interested in investigating the consequences of these structural
features, independently of a specific physical realization.
Within this line of research, the nonlinear heat flow constrained to the $L^2$-unit
sphere of a Hilbert space has emerged as a natural and analytically tractable model.
The projection operator $\pi$ ensures that the modified vector field
$\widetilde{f}(x)=\pi(x)[f(x)]$ is tangent to the manifold, thus preserving invariance.
This approach, initiated by Zdzislaw and Javed
\cite{Brz+Huss_2024,Hussain_2005,Huss_2023},
investigated that the projected nonlinear heat equation admits global solutions,
generates a gradient flow on the manifold, and maintains invariance of the constraint.

The projection of parabolic dynamics onto constrained manifolds has been examined in
several seminal works. Rybka \cite{Rybka} and Caffarelli--Lin \cite{Cafe} considered
the heat equation in $\mathcal{L}^{2}(\mathcal{O})$ under algebraic constraints of the
form
\begin{equation*}
\mathcal{M} = \big\{ u \in \mathcal{L}^{2}(\mathcal{O}) \cap C(\mathcal{O})
: \int_{\mathcal{O}} u^{k}(x)  dx = C_{k},\;  k = 1, 2, \dots, N \big\},
\end{equation*}
where $\mathcal{O}$ is a bounded, connected region in $\mathbb{R}^{2}$.
Rybka established the global existence and uniqueness of solutions to the projected
heat flow
\begin{equation}
\begin{gathered}
\frac{d u}{d t} = \Delta u - \sum_{k=1}^{N} \lambda_{k} u^{k-1} \quad
 \text{in } \mathcal{O} \subset \mathbb{R}^{2}, \\
\frac{\partial u}{\partial n} = 0 \quad \text{on } \partial \mathcal{O}, \\
u(0,x) = u_{0}
\end{gathered} \label{eq1}
\end{equation}
with Lagrange multipliers $\lambda_{k}(u)$ chosen so that $u_{t}$ is orthogonal
to $\mathrm{Span}\{u^{k-1}\}$. He proved that solutions converge asymptotically to
constant states. In related work, Caffarelli and Lin \cite{Cafe} proved global
well-posedness of energy-conserving heat flows, and extended their framework to
singularly perturbed nonlocal parabolic systems, establishing convergence to weak
solutions of limiting constrained flows.

The constrained dynamics paradigm has since been extended to fluid equations and
other PDEs. For example, Brze\'zniak, Dhariwal, and Mariani \cite{Gaurav}
studied two-dimensional Navier--Stokes equations under an energy constraint,
proving global existence, uniqueness, and convergence to Euler flows as viscosity
tends to zero. These works illustrate the broader significance of projecting
dissipative PDEs onto constraint manifolds.

This article continues the line of work on parabolic evolutions constrained to manifolds
initiated for heat-type flows in \cite{Brz+Huss_2024,Cafe,Hussain_2005,Huss_2023,Rybka}
and subsequently developed for other dissipative equations. In contrast with the linear
heat flow under algebraic constraints treated in \cite{Cafe,Rybka}, we consider a
power-type reaction term under the $\mathcal{L}^{2}$-unit-sphere constraint and
homogeneous Dirichlet boundary conditions. This leads to a constrained evolution driven
by the Dirichlet Laplacian and a superlinear Nemytskii nonlinearity, where the constraint
is enforced through the orthogonal projection onto the tangent bundle of $\mathcal{M}$.
Within this framework, the contribution is threefold. First, we derive an explicit
projected vector field and verify that the induced semiflow is the gradient flow of the
energy restricted to $\mathcal{M}$, so that the energy identity provides a Lyapunov
structure on the manifold. Second, we establish a constrained \L{o}jasiewicz-Simon
inequality near every stationary point $\varphi\in\mathcal{S}$ in a form adapted to the
Riemannian structure of the $\mathcal{L}^{2}$-sphere; this yields convergence of
every global trajectory to a single equilibrium and permits a sharp dichotomy between
exponential decay near nondegenerate equilibria and algebraic decay otherwise.
Third, combining the resulting asymptotic compactness with the dissipative bounds
supplied by the Lyapunov functional, we prove the existence of a compact global
attractor in $\mathcal{V}\cap\mathcal{M}$.

The broader context of convergence to equilibrium in parabolic flows is well
studied (see \cite{BH_LDHilbert_2026,Brz+Huss_2020,Hale,Haraux,Haraux2011,
Haraux2009,Vishnevski,Lions,Matano,Simon}).
For decay rates we refer to \cite{Haraux2003,Songu}, and for attractor theory in
dissipative PDEs to \cite{Carvalho,Songu,Kapustyan,Temam}. Our results place the
projected nonlinear heat equation firmly within this classical framework while
highlighting its structural and geometric features.

\paragraph{Outline of the paper.}
Section~2 collects notations, functional spaces, and preliminary results.
Section~3 establishes pre-compactness of solution orbits, characterizes the
$\omega$-limit set, and proves a \L{o}jasiewicz-Simon inequality on the manifold.
Section~4 derives decay rates for convergence to equilibria, distinguishing
exponential and algebraic regimes. Section~5 proves the existence of a global
attractor in $\mathcal{M}\cap H^1$. Together, these results provide a mathematically
rigorous and comprehensive description of the asymptotic dynamics of the constrained
nonlinear heat flow.

\section{Functional settings, assumptions and preliminaries}
Let us set the notation that will be used throughout this paper.

\subsection{Manifold and projection}

In this paper we work with the unit sphere
\[
\mathcal{M} = \{ u \in \mathcal{H} : |u|_{\mathcal{H}}^{2} = 1\}
\]
as a smooth submanifold of the Hilbert space $\mathcal{H}$ endowed with inner product
$\langle\cdot,\cdot\rangle$. It is well known that $\mathcal{M}$ is a Hilbert
(Riemannian) submanifold of $\mathcal{H}$. For any $a \in \mathcal{M}$, the tangent
space is
\[
T_{a}\mathcal{M} = \{ v \in \mathcal{H} : \langle a, v \rangle = 0 \}.
\]
Let $\pi_{a} : \mathcal{H} \longrightarrow T_{a}\mathcal{M}$
be the orthogonal projection onto $T_a\mathcal{M}$.

\begin{lemma}
\label{lem-orthogonal projection}
If $a \in \mathcal{M}$, then for all $v \in \mathcal{H}$,
\[
\pi_{a}(v) = v - \langle a, v \rangle a.
\]
\end{lemma}

\begin{remark}\label{rem-1-spaces} \rm
Let $\mathcal{O} \subset \mathbb{R}^{d}$ be a bounded $C^{2}$ domain and let
$n \in [1, \infty)$ be such that the Sobolev embedding
\[
\mathcal{H}^{1,2}(\mathcal{O}) \hookrightarrow \mathcal{L}^{2n}(\mathcal{O})
\]
holds. Equivalently:
\begin{itemize}
\item if $d = 2$, the embedding holds for every finite $2n$;
\item if $d \ge 3$, we assume $2n \le 2^{\ast} := \frac{2d}{d-2}$, i.e.
\begin{equation}
\frac{1}{d} \geq \frac{1}{2} - \frac{1}{2n}.
\label{m_thm}
\end{equation}
\end{itemize}
Throughout we adopt the following spaces:
\[
\mathcal{H} = \mathcal{L}^{2}(\mathcal{O}), \quad
\mathcal{V} = \mathcal{H}_{0}^{1,2}(\mathcal{O}), \quad
\mathcal{E} = \mathcal{D}(A),
\]
where $A$ is the (negative) Laplace operator with homogeneous Dirichlet boundary
conditions,
\begin{equation} \label{1-1.1}
\begin{gathered}
\mathcal{D}(A) = \mathcal{H}_{0}^{1,2}(\mathcal{O})
\cap \mathcal{H}^{2,2}(\mathcal{O}),  \\
Au = -\Delta u, \quad u \in \mathcal{D}(A).
\end{gathered}
\end{equation}
It is classical (see \cite[Theorem 4.1.2, p.~79]{I.Vrabie}) that $A$ is
self-adjoint and positive on $\mathcal{H}$, that $\mathcal{V} = \mathcal{D}(A^{1/2})$,
and
\[
\|u\|^{2} = |A^{1/2}u|_{\mathcal{H}}^{2} = \int_{\mathcal{O}} |\nabla u(x)|^{2}  dx.
\]
Moreover,
\[
\mathcal{E} \subset \mathcal{V} \subset \mathcal{H} \subset \mathcal{V}'
=: \mathcal{H}^{-1}(\mathcal{O}),
\]
with all injections continuous and dense.
\end{remark}

For each $T \ge 0$ we set
\begin{equation}
\mathcal{X}_{T} := \mathcal{L}^{2}(0,T;\mathcal{E}) \cap C([0,T];\mathcal{V}).
\label{eqn-X_T}
\end{equation}
Then $\mathcal{X}_{T}$ is a Banach space with norm
\begin{equation}
|u|_{\mathcal{X}_{T}}^{2}
= \sup_{t \in [0,T]} \|u(t)\|^{2} + \int_{0}^{T} |u(t)|_{\mathcal{E}}^{2}  dt, \quad u \in \mathcal{X}_{T}.
\label{eqn-X_T-norm}
\end{equation}

\begin{corollary}
\label{cor-projection}
Within the framework of Remark \ref{rem-1-spaces}, for any
$u \in \mathcal{E} \cap \mathcal{M}$ we have the identity
\begin{equation}
\pi_{u} \left( \Delta u - |u|^{2n-2} u \right)
= \Delta u - |u|^{2n-2} u + \left( \|u\|^{2} + |u|_{\mathcal{L}^{2n}}^{2n} \right) u.
\label{1-1.2}
\end{equation}
\end{corollary}

\begin{proof}
Let $u \in \mathcal{E} \cap \mathcal{M}$. Then $\Delta u, |u|^{2n-2}u \in \mathcal{H}$
by elliptic regularity and the embedding from Remark \ref{rem-1-spaces}.
By Lemma \ref{lem-orthogonal projection},
\begin{equation} \label{1-1.3}
\begin{aligned}
\pi_{u} \left( \Delta u - |u|^{2n-2} u \right)
&= \Delta u - |u|^{2n-2} u - \langle \Delta u - |u|^{2n-2} u, u \rangle u  \\
&= \Delta u - |u|^{2n-2} u + \langle \nabla u, \nabla u \rangle u
 + \langle |u|^{2n-2} u, u \rangle u  \\
&= \Delta u - |u|^{2n-2} u + \left( \|u\|^{2} + |u|_{\mathcal{L}^{2n}}^{2n} \right) u.
\end{aligned}
\end{equation}
Here we used integration by parts with Dirichlet boundary conditions,
cf. \cite[Corollary 8.10, p.~82]{Haim}, and the identity
$\int_{\mathcal{O}} |u|^{2n-2} u \cdot u = \|u\|_{\mathcal{L}^{2n}}^{2n}$.
\end{proof}

The next result summarizes the well-posedness and gradient-flow structure for the
projected equation; the ensuing energy relation will be used repeatedly.

\begin{theorem}[\cite{Hussain_2005}]
\label{1-Global}
For every $u_{0} \in \mathcal{H}_{0}^{1,2}(\mathcal{O}) \cap \mathcal{M}$, where
$\mathcal{M} = \{ u \in \mathcal{L}^{2}(\mathcal{O}) : |u|_{\mathcal{L}^{2}(\mathcal{O})}
 = 1 \}$,
there exists a unique function $u : [0, \infty) \to \mathcal{H}_{0}^{1,2}(\mathcal{O})$
such that $u \in \mathcal{X}_{T}$ for every $T > 0$ and
\begin{equation}\label{eqn-main}
\begin{gathered}
\frac{\partial u}{\partial t} = \Delta u - |u|^{2n-2} u + \left( \|u\|^{2} + |u|_{\mathcal{L}^{2n}}^{2n} \right) u, \\
u(0) = u_{0}.
\end{gathered}
\end{equation}
Moreover, $u(t) \in \mathcal{M}$ for all $t \ge 0$. In particular, the energy
identity holds:
\begin{equation}
\Phi(u(t)) - \Phi(u_{0}) = -\int_{0}^{t} | \frac{du}{dt}(s) |_{\mathcal{H}}^{2}  ds,
\quad t \ge 0,
\label{1-energy_equality}
\end{equation}
and
\[
\|u(t)\| \le 2 \Phi(u_{0}), \quad t \ge 0,
\]
where $\Phi : \mathcal{V} \to \mathbb{R}$ is defined by
\[
\Phi(u) = \frac{1}{2} |\nabla u|_{\mathcal{L}^2(\mathcal{O})}^{2}
+ \frac{1}{2n} |u|_{\mathcal{L}^{2n}(\mathcal{O})}^{2n}, \quad n \in \mathbb{N}.
\]
\end{theorem}

\begin{remark} \rm
The existence and uniqueness asserted in Theorem~\ref{1-Global} are obtained by
a standard contraction argument applied to the mild formulation of \eqref{eqn-main}.
The Dirichlet Laplacian generates an analytic semigroup on
$\mathcal{L}^{2}(\mathcal{O})$,
and under the standing restrictions on $n$ and the spatial dimension, the Nemytskii map
$u\mapsto |u|^{2n-2}u$ is locally Lipschitz from $\mathcal{H}_{0}^{1,2}(\mathcal{O})$
into $\mathcal{L}^{2}(\mathcal{O})$.
The additional lower-order term involving $\|u\|^{2}$ and $|u|_{\mathcal{L}^{2n}}^{2n}$
is smooth on bounded sets and does not affect the contraction property on short time intervals.
The invariance of $\mathcal{M}$ follows from the fact that the right-hand side of
\eqref{eqn-main} is orthogonal in $\mathcal{L}^{2}(\mathcal{O})$ to $u(t)$, which implies
$\frac{d}{dt}|u(t)|_{\mathcal{L}^{2}}^{2}=0$.
Global existence is then obtained by combining the local theory with the energy identity
\eqref{1-energy_equality}, which prevents finite-time blow-up.
\end{remark}

For convenience we recall some definitions and abstract results on convergence
to equilibrium and global attractors from Chapter~6 of \cite{Songu}, which will be
used throughout.

\begin{corollary}[\cite{Songu}]  \label{2.2}
Suppose that $\mathcal{H}$ is a complete metric space and $S(t)$ is a nonlinear
$C_0$-semigroup on $\mathcal{H}$. Let $x \in \mathcal{H}$. If there exists $t_0 \ge 0$
such that
\[
\{ u(t) : t > t_0 \} := \cup_{t > t_0} S(t) x
\]
is relatively compact in $\mathcal{H}$, then the $\omega$-limit set $\omega(x)$ is a
compact, connected, invariant set.
\end{corollary}

\begin{theorem}\cite{Songu}
Let $\Gamma : \mathbb{R}^N \to \mathbb{R}$ be analytic in a neighborhood of a point
$a \in \mathbb{R}^N$. Then there exist $\sigma > 0$ and $\theta \in (0, \tfrac12]$
such that for all $x \in \mathbb{R}^N$,
\[
\|x - a\| < \sigma \implies  \|\nabla \Gamma(x)\| \ge |\Gamma(x)
- \Gamma(a)|^{1 - \theta}.
\]
\end{theorem}

\begin{definition}[\cite{Songu}] \rm
Suppose that $\mathcal{H}$ is a complete metric space, and $S(t)$ is a nonlinear
 $C_0$-semigroup on $\mathcal{H}$. A set $\mathcal{A} \subset \mathcal{H}$ is an
 \emph{attractor} if:
\begin{itemize}
\item[(i)] $\mathcal{A}$ is invariant: $S(t)\mathcal{A} = \mathcal{A}$ for all
$t \ge 0$;
\item[(ii)] there exists an open neighborhood $\mathcal{U}$ of $\mathcal{A}$ such
that for every $u_0 \in \mathcal{U}$,
\[
\operatorname{dist} \left( S(t) u_0, \mathcal{A} \right) = \inf_{y \in \mathcal{A}}
d \left( S(t) u_0, y \right) \longrightarrow 0 \quad \text{as } t \to \infty.
\]
\end{itemize}
\end{definition}

\begin{definition}[\cite{Songu}] \rm
If $\mathcal{A}$ is a compact attractor and it attracts all bounded subsets of
$\mathcal{H}$, then $\mathcal{A}$ is called a \emph{global} (or \emph{universal})
attractor.
\end{definition}

\begin{theorem}[\cite{Songu}] \label{attractor}
Suppose that $\mathcal{H}$ is a Banach space and $S(t)$ is a nonlinear $C_0$-semigroup
on $\mathcal{H}$ such that:
\begin{itemize}
\item[(i)] there exists a bounded absorbing set $\mathcal{B}_0$;
\item[(ii)] for any bounded set $\mathcal{B}$ there exists $t_0(\mathcal{B}) \ge 0$
with
$\cup_{t \ge t_0(\mathcal{B})} S(t) \mathcal{B}$
relatively compact in $\mathcal{H}$.
\end{itemize}
Then $\mathcal{A} = \omega(\mathcal{B}_0)$ is a global attractor.
\end{theorem}

We finally recall the following compactness lemma
(see \cite[Chapter~3, Lemma~1.2]{Temam}).

\begin{lemma}[{\cite[Lemma III 1.2]{Temam}}] \label{1-tem}
Let $\mathcal{V}, \mathcal{H}, \mathcal{V}'$ be Hilbert spaces with $\mathcal{V}'$
the dual of $\mathcal{V}$ and
\[
\mathcal{V} \hookrightarrow \mathcal{H} \cong \mathcal{H}'
\hookrightarrow \mathcal{V}',
\]
with dense continuous embeddings. If $u \in \mathcal{L}^{2}(0,T;\mathcal{V})$ and its
weak derivative $u_t \in \mathcal{L}^{2}(0,T;\mathcal{V}')$, then there exists
$\tilde{u} \in \mathcal{L}^{2}(0,T;\mathcal{V}) \cap C([0,T];\mathcal{V})$ such that
$\tilde{u} = u$ a.e., and
\[
| u(t) |^{2} = | u_{0} |^{2} + 2 \int_{0}^{t} \langle u'(s), u(s) \rangle  ds, \quad
\forall t \in [0,T].
\]
\end{lemma}

\section{\L{o}jasiewicz-Simon inequality and convergence to equilibrium}

In this section we prove convergence to equilibrium for the global solutions constructed
in Theorem~\ref{1-Global}. The argument combines two ingredients. On the one hand, the
orbit $\{u(t):t\ge 1\}$ is shown to be relatively compact in $\mathcal{V}$, which
implies that the $\omega$-limit set is a nonempty compact subset of $\mathcal{V}$.
On the other hand, near each stationary point $\varphi\in\mathcal{S}$ we establish
a \L{o}jasiewicz-Simon inequality for the energy restricted to the manifold
$\mathcal{M}$, formulated in terms of the projected gradient
$\pi_u(\Delta u-|u|^{2n-2}u)$. The convergence theorem then follows from the energy
identity \eqref{1-energy_equality} and the fact that the orbit ultimately remains in
a neighborhood where a single \L{o}jasiewicz-Simon inequality applies.


\begin{theorem}\label{thm:precompact}
Let $u_0\in \mathcal V\cap\mathcal M$, and let $u(\cdot)$ be the global solution of
\eqref{eqn-main} given by Theorem~\ref{1-Global}. Then the orbit $\{u(t): t\ge 0\}$
is relatively compact in $\mathcal V$. Equivalently, for every sequence $t_k\to\infty$
there exist a subsequence (not relabelled) and $u^\ast\in\mathcal V$ such that
$u(t_k)\to u^\ast$ in $\mathcal V$.
\end{theorem}

\begin{proof}
 For $\alpha\in(0,1)$ denote by $A^\alpha$ the fractional
powers of $A$ and set $\mathcal D(A^\alpha)$ as usual; the compactness of the resolvent
implies that the embedding $\mathcal D(A^\alpha)\hookrightarrow \mathcal H$ is
compact for every $\alpha>0$.
Moreover, for $\alpha>\frac12$ one has the continuous embedding
\begin{equation}\label{eq:Dalpha-to-V}
\mathcal D(A^\alpha)\hookrightarrow \mathcal V=H_0^1(\mathcal O),
\end{equation}
and the embedding $\mathcal D(A^\alpha)\hookrightarrow \mathcal V$ is compact.
This follows, for
instance, from the spectral representation of $A$ and the characterisation
$\mathcal D(A^{1/2})=H_0^1(\mathcal O)$.

Write \eqref{eqn-main} in the semilinear form
\begin{equation}\label{eq:semilinear-form}
u'(t)+A u(t)=F(u(t)),\quad t>0,
\end{equation}
where
\[
F(u):= -|u|^{2n-2}u + \bigl(\|u\|_{\mathcal V}^{2}+|u|_{L^{2n}}^{2n}\bigr)u.
\]
Fix $\alpha\in(\tfrac12,1)$ and $T>0$. The variation-of-constants formula for
\eqref{eq:semilinear-form} gives, for all $t\ge T$,
\begin{equation}\label{eq:variation-constants-shift}
u(t)=e^{-(t-T)A}u(T)+\int_T^{t} e^{-(t-s)A}F(u(s))\,ds.
\end{equation}
Applying $A^\alpha$ and using analyticity of the semigroup yields
\begin{equation}\label{eq:smoothing-est}
\|A^\alpha e^{-\tau A}\|_{\mathcal L(\mathcal H)}
\le C_\alpha\,\tau^{-\alpha},\quad \tau\in(0,1],
\end{equation}
and $\|A^\alpha e^{-\tau A}\|_{\mathcal L(\mathcal H)}\le C_\alpha$ for $\tau\ge 1$.

The energy identity \eqref{1-energy_equality} implies that $t\mapsto \Phi(u(t))$ is
non-increasing,
hence $\sup_{t\ge 0}\Phi(u(t))\le \Phi(u_0)$. In particular,
\begin{equation}\label{eq:Vbound}
\sup_{t\ge 0}\|u(t)\|_{\mathcal V}<\infty,
\end{equation}
since $\Phi(u)\ge \frac12\|\nabla u\|_{L^2}^2$. The constraint $u(t)\in\mathcal M$ gives
$|u(t)|_{\mathcal H}=1$ for all $t\ge 0$. Under the standing restriction on $n$ and the
spatial dimension ensuring $\mathcal V\hookrightarrow L^{2n}(\mathcal O)$,
the bound \eqref{eq:Vbound} yields
\begin{equation}\label{eq:L2n-bound}
\sup_{t\ge 0} |u(t)|_{L^{2n}}<\infty.
\end{equation}
Consequently, each term in $F(u(t))$ is uniformly bounded in $\mathcal H$ for $t\ge 0$.
Indeed, the Nemytskii map $u\mapsto |u|^{2n-2}u$ maps $L^{2n}(\mathcal O)$ into
$L^2(\mathcal O)$, and \eqref{eq:L2n-bound} shows that
$|\,|u(t)|^{2n-2}u(t)|_{\mathcal H}$ is bounded uniformly in
time. The remaining lower-order term is bounded in $\mathcal H$ as well because
$\|u(t)\|_{\mathcal V}$ and $|u(t)|_{L^{2n}}$ are uniformly bounded and
$|u(t)|_{\mathcal H}=1$.
Hence there exists $M>0$ such that
\begin{equation}\label{eq:Fbound}
\sup_{t\ge 0}\,|F(u(t))|_{\mathcal H}\le M.
\end{equation}

Fix any sequence $t_k\to\infty$. Choose $T\ge 1$ so that $t_k\ge T+1$ for all $k$ large,
and write
\eqref{eq:variation-constants-shift} at $t=t_k$:
\begin{equation}\label{eq:decomp-tk}
u(t_k)=e^{-(t_k-T)A}u(T)+\int_T^{t_k} e^{-(t_k-s)A}F(u(s))\,ds.
\end{equation}
Applying $A^\alpha$ and estimating with \eqref{eq:smoothing-est} and
\eqref{eq:Fbound} yields
\begin{align*}
\|A^\alpha u(t_k)\|_{\mathcal H}
&\le \|A^\alpha e^{-(t_k-T)A}u(T)\|_{\mathcal H}
     +\int_T^{t_k}\|A^\alpha e^{-(t_k-s)A}F(u(s))\|_{\mathcal H}\,ds\\
&\le C_\alpha\|u(T)\|_{\mathcal H}+ C_\alpha M\int_0^{t_k-T} \tau^{-\alpha}\,d\tau.
\end{align*}
Since $\alpha<1$, the integral is finite and bounded uniformly in $k$:
\[
\int_0^{t_k-T}\tau^{-\alpha}\,d\tau
\le \int_0^{\infty}\tau^{-\alpha}\mathbf 1_{(0,1]}(\tau)\,d\tau
+\int_1^\infty \tau^{-\alpha}\,d\tau <\infty.
\]
Thus $\{u(t_k)\}$ is bounded in $\mathcal D(A^\alpha)$.

Because $A$ has compact resolvent, the embedding
$\mathcal D(A^\alpha)\hookrightarrow \mathcal V$
is compact for every $\alpha>\frac12$ by \eqref{eq:Dalpha-to-V}.
Therefore, the bounded sequence
$\{u(t_k)\}$ admits a subsequence converging in $\mathcal V$.
This proves the relative compactness
of the orbit in $\mathcal V$.
\end{proof}

\begin{corollary}\label{omega}
The $\omega$-limit set
$\omega(u_0)=\bigcap_{r\ge1}\overline{\{u(t):t\ge r\}}^{\mathcal{V}}$ exists and
is compact in $\mathcal{V}$.
\end{corollary}

\begin{proof}
From last corollary, $\left\{ u(t):t\geq r\right\} $ is pre-compact in $\mathcal{V}$
for all $r\geq 1$. Since closure of pre-compact is also pre-compact so
$\overline{\left\{ u(t):t\geq r\right\} }$ is also pre-compact for all
 $r\geq1 $. Further $\overline{\left\{ u(t):t\geq r\right\} }$ is closed and hence
complete in $\mathcal{V}$ norm therefore $\overline{\left\{ u(t):t\geq r\right\} }$
being pre-compact and complete, it follows that $\overline{\left\{ u(t):t\geq r\right\} }$ is compact for all $r\geq 1$, in $\mathcal{V}$. Thus $\omega (u_{0})$ being decreasing intersection of
non-empty compact sets in $\mathcal{V}$, is non-empty and compact in $\mathcal{V}$.
\end{proof}

\subsection*{\L{o}jasiewicz-Simon inequality on $\mathcal{M}$}
Let $\mathcal{S}$ denote the set of stationary points of \eqref{eqn-main} on
$\mathcal{M}$ (critical points of $\Phi$ under the constraint $|u|_{\mathcal{H}}=1$).
The next result is a constrained \L{o}jasiewicz-Simon inequality; the proof follows
the abstract LS theory (see, e.g., \cite{Songu} and the references therein) and the
argument of Jendoubi \cite{Jendoubi}, adapted to the present manifold setting.


\begin{theorem}[\L{o}jasiewicz-Simon inequality]\label{Lwoj}
Let $\varphi\in \mathcal V\cap\mathcal M$ be a stationary point of \eqref{eqn-main}.
Then there exist constants $\sigma>0$, $\theta\in(0,\tfrac12]$, and $C>0$ such that
\begin{equation}\label{eq:LS}
|\Phi(u)-\Phi(\varphi)|^{1-\theta}\le C \|\nabla_{\mathcal M}\Phi(u)\|_{\mathcal H},
\quad
u\in \mathcal V\cap\mathcal M, \|u-\varphi\|_{\mathcal V}<\sigma,
\end{equation}
where $\nabla_{\mathcal M}\Phi(u)$ denotes the $\mathcal H$-gradient of $\Phi$
restricted to the manifold $\mathcal M$.
\end{theorem}

\begin{proof}
The manifold is
\[
\mathcal M=\{u\in \mathcal H:\ G(u)=0\},\quad G(u)
:=\frac12\bigl(|u|_{\mathcal H}^{2}-1\bigr),
\]
and $G$ is a real-analytic mapping $\mathcal H\to\mathbb R$. Its derivative satisfies
$G'(u)h=(u,h)_{\mathcal H}$, hence $G'(\varphi)\neq 0$ and the regular value theorem
yields that
$\mathcal M$ is a real-analytic embedded submanifold of $\mathcal H$ in a
neighbourhood of
$\varphi$. The tangent space is
\[
T_\varphi\mathcal M=\ker G'(\varphi)=\{h\in\mathcal H:\ (h,\varphi)_{\mathcal H}=0\}.
\]
Write $P_\varphi:\mathcal H\to T_\varphi\mathcal M$ for the orthogonal projection
$P_\varphi h=h-(h,\varphi)_{\mathcal H}\varphi$.

A convenient analytic chart is obtained by the normalisation map. Let
\[
U:=\{h\in T_\varphi\mathcal M:\ |h|_{\mathcal H}<\tfrac12\},
\quad
\Psi(h):=\frac{\varphi+h}{|\varphi+h|_{\mathcal H}},\quad h\in U.
\]
Since $|\varphi|_{\mathcal H}=1$ and $h\perp \varphi$, one has
$|\varphi+h|_{\mathcal H}^{2}=1+|h|_{\mathcal H}^{2}$, so the denominator is bounded
away from
$0$ on $U$. The map $h\mapsto (1+|h|_{\mathcal H}^{2})^{-1/2}$ is real-analytic on
$\{|h|_{\mathcal H}<1\}$, and therefore $\Psi:U\to\mathcal H$ is real-analytic.
Moreover,
$\Psi(U)\subset\mathcal M$ by construction, $\Psi(0)=\varphi$, and the differential
satisfies
$D\Psi(0)=\mathrm{Id}$ on $T_\varphi\mathcal M$. In particular, $\Psi$ is a
real-analytic immersion.

To see that $\Psi$ parametrises a full neighbourhood of $\varphi$ in $\mathcal M$,
consider the analytic map
\[
\Xi:\mathcal M\cap B_{\mathcal H}(\varphi,\tfrac12)\to T_\varphi\mathcal M,
\quad
\Xi(u):=P_\varphi(u-\varphi).
\]
If $u\in\mathcal M$, then $|u|_{\mathcal H}=|\varphi|_{\mathcal H}=1$ and the identity
$|u-\varphi|_{\mathcal H}^{2}=2\bigl(1-(u,\varphi)_{\mathcal H}\bigr)$ implies that
$(u,\varphi)_{\mathcal H}>0$ whenever $u$ is sufficiently close to $\varphi$. For such $u$ set
$\alpha(u):=(u,\varphi)_{\mathcal H}\in(0,1]$ and $h:=\Xi(u)\in T_\varphi\mathcal M$.
Then
$u=\alpha(u)\varphi+h$ with $h\perp\varphi$ and
$|u|_{\mathcal H}^{2}=\alpha(u)^{2}+|h|_{\mathcal H}^{2}=1$,
so $\alpha(u)=\sqrt{1-|h|_{\mathcal H}^{2}}$. Hence
\[
u=\sqrt{1-|h|_{\mathcal H}^{2}}\,\varphi+h
=\frac{\varphi+\frac{h}{\sqrt{1-|h|_{\mathcal H}^{2}}}}
{\big|\varphi+\frac{h}{\sqrt{1-|h|_{\mathcal H}^{2}}}\big|_{\mathcal H}}
=\Psi\Big(\frac{h}{\sqrt{1-|h|_{\mathcal H}^{2}}}\Big).
\]
For $u$ close to $\varphi$, $|h|_{\mathcal H}$ is small and
$h\mapsto h/\sqrt{1-|h|_{\mathcal H}^{2}}$ maps a neighbourhood of $0$ in
$T_\varphi\mathcal M$
into $U$ real-analytically. This shows that $\Psi$ is a real-analytic bijection from a
neighbourhood of $0$ in $T_\varphi\mathcal M$ onto a neighbourhood of $\varphi$
in $\mathcal M$,
with analytic inverse, hence a real-analytic chart.

We define the reduced functional on the Hilbert space $T_\varphi\mathcal M$ by
\[
\mathcal F(h):=\Phi(\Psi(h)),\quad h\in U.
\]
Since $\Phi$ is real-analytic on $\mathcal V$ and $\Psi(h)\in\mathcal V$ for $h$
small (because
$\Psi(h)$ differs from $\varphi$ by a small $\mathcal H$--perturbation and
$\varphi\in\mathcal V$),
the map $\mathcal F$ is real-analytic on $U$ as a function on $T_\varphi\mathcal M$.
The point $h=0$ is critical: by the chain rule,
\[
D\mathcal F(0)k
=\langle \Phi'(\varphi),D\Psi(0)k\rangle
=\langle \Phi'(\varphi),k\rangle,
\quad k\in T_\varphi\mathcal M,
\]
and stationarity of $\varphi$ means precisely that $\Phi'(\varphi)$ annihilates
$T_\varphi\mathcal M$, so $D\mathcal F(0)=0$.

Let $\nabla\mathcal F(h)$ denote the gradient of $\mathcal F$ in the Hilbert space
$T_\varphi\mathcal M$ with the inherited $\mathcal H$ inner product.
For $k\in T_\varphi\mathcal M$,
\[
(\nabla\mathcal F(h),k)_{\mathcal H}=D\mathcal F(h)k
=\bigl(\nabla_{\mathcal H}\Phi(\Psi(h)),D\Psi(h)k\bigr)_{\mathcal H}.
\]
Set $u=\Psi(h)$. Since $\Psi$ takes values in $\mathcal M$, one has
$\mathrm{Ran}\,D\Psi(h)\subset T_u\mathcal M$.
Write $\pi_u$ for the orthogonal projection $\mathcal H\to T_u\mathcal M$. Then
\[
\bigl(\nabla_{\mathcal H}\Phi(u),D\Psi(h)k\bigr)_{\mathcal H}
=\bigl(\pi_u\nabla_{\mathcal H}\Phi(u),D\Psi(h)k\bigr)_{\mathcal H}.
\]
The operator $D\Psi(h):T_\varphi\mathcal M\to T_u\mathcal M$ is a bounded isomorphism
for $h$ sufficiently small, because $D\Psi(0)=\mathrm{Id}$ and $D\Psi$ depends
continuously on $h$. Consequently the operator
\[
\mathcal A(h):=\bigl(D\Psi(h)\bigr)^{\!*}:T_u\mathcal M\to T_\varphi\mathcal M
\]
is bounded and invertible for $h$ small, and
\begin{equation}\label{eq:grad-relation}
\nabla\mathcal F(h)=\mathcal A(h)\bigl(\pi_u\nabla_{\mathcal H}\Phi(u)\bigr).
\end{equation}
In particular, since $\mathcal A(h)$ varies continuously and is invertible near $0$,
there exist $c_1,c_2>0$ and a neighbourhood $U_0\subset U$ of $0$ such that for all
$h\in U_0$,
\begin{equation}\label{eq:norm-equivalence}
c_1\|\pi_{\Psi(h)}\nabla_{\mathcal H}\Phi(\Psi(h))\|_{\mathcal H}
\le \|\nabla\mathcal F(h)\|_{\mathcal H}
\le c_2\|\pi_{\Psi(h)}\nabla_{\mathcal H}\Phi(\Psi(h))\|_{\mathcal H}.
\end{equation}

Let $L:=D(\nabla\mathcal F)(0)$ be the linearization of the gradient map at $0$. A direct
differentiation of \eqref{eq:grad-relation} at $h=0$, using $D\Psi(0)=\mathrm{Id}$ and
the fact that $\varphi$ is a critical point, shows that $L$ is the restriction of the
second variation of $\Phi$ to $T_\varphi\mathcal M$. In the present setting, $L$ is a
self-adjoint operator on $T_\varphi\mathcal M$ with compact resolvent, inherited from
the elliptic part of $\Phi$ through
the Dirichlet Laplacian. This is the spectral hypothesis required in the
\L{o}jasiewicz-Simon theorem \cite{Simon}.

The classical \L{o}jasiewicz-Simon inequality applied to the real-analytic functional
$\mathcal F$ at the critical point $0$ therefore yields constants $\rho>0$,
$\theta\in(0,\tfrac12]$, and $C_0>0$ such that
\begin{equation}\label{eq:LS-chart}
|\mathcal F(h)-\mathcal F(0)|^{1-\theta}\le C_0\|\nabla\mathcal F(h)\|_{\mathcal H},
\quad h\in T_\varphi\mathcal M,\ |h|_{\mathcal H}<\rho.
\end{equation}
Returning to $u=\Psi(h)$, one has $\mathcal F(h)=\Phi(u)$ and
 $\mathcal F(0)=\Phi(\varphi)$.
Combining \eqref{eq:LS-chart} with the right-hand inequality in
\eqref{eq:norm-equivalence} gives
\[
|\Phi(u)-\Phi(\varphi)|^{1-\theta}
\le C_0\|\nabla\mathcal F(h)\|_{\mathcal H}
\le C_0 c_2\,\|\pi_u\nabla_{\mathcal H}\Phi(u)\|_{\mathcal H}.
\]
Since $\Psi$ is a local diffeomorphism between neighbourhoods of $0$ and $\varphi$,
smallness of $h$ is equivalent to smallness of $u-\varphi$ in $\mathcal H$, and hence,
on bounded sets, also in $\mathcal V$ by elliptic regularity for the resolvent of $A$.
Shrinking the neighbourhood if necessary yields \eqref{eq:LS} with $\sigma>0$ and
$C:=C_0c_2$.
\end{proof}

\begin{thebibliography}{99}
\bibitem[1]{BH_LDHilbert_2026} Z. Brze\'zniak, J. Hussain; \emph{Large deviations for the stochastic nonlinear heat equation on a Hilbert manifold}, Potential Analysis, \textbf{64}, 36, 2026.
\bibitem[2]{Brz+Huss_2020} Z. Brze\'zniak, J. Hussain; \emph{Global solution of nonlinear stochastic heat equation with solutions in a Hilbert manifold}, Stochastics and Dynamics, \textbf{20} (2020), 6, 2040012.
\bibitem[3]{Brz+Huss_2024} Z. Brze\'zniak, J. Hussain; \emph{Global solution of nonlinear heat equation with solutions in a Hilbert manifold}, Nonlinear Analysis, \textbf{242}  (2024), 113505.
\bibitem[5]{Haim} H. Brezis; \emph{Functional Analysis, Sobolev Spaces and Partial Differential Equations}, Spring, 2010.
\bibitem[6]{Gaurav} Z. Brze\'{z}niak, G. Dhariwal, M. Mariani; \emph{2D constrained Navier--Stokes equations}, Journal of Differential Equations, \textbf{263}  (2017), 12, 8282--8327.
\bibitem[7]{Cafe} L. Caffarelli, F. Lin; \emph{Nonlocal heat flows preserving the $\mathcal{L}^{2}$-energy}, Discrete and Continuous Dynamical Systems, \textbf{23} (2009), 1\&2, 49--64.
\bibitem[8]{Carvalho} A. N. de Carvalho, J. W. Cholewa, T. Dlotko; \emph{Examples of global attractors in parabolic problems}, Hokkaido Mathematical Journal, \textbf{27}  (1998), 1, 77--103.
\bibitem[12]{Hale} J. K. Hale, G. Raugel; \emph{Convergence in gradient-like systems with applications to PDE}, Zeitschrift f\"{u}r Angewandte Mathematik und Physik, \textbf{43} (1992), 1, 63--124.
\bibitem[13]{Haraux} A. Haraux, M. A. Jendoubi; \emph{Convergence of bounded weak solutions of the wave equation with dissipation and analytic nonlinearity}, Calculus of Variations and Partial Differential Equations, \textbf{9} (1999), 2, 95--124.
\bibitem[14]{Haraux2003} A. Haraux, M. A. Jendoubi, O. Kavian; \emph{Rate of decay to equilibrium in some semilinear parabolic equations}, Journal of Evolution Equations, \textbf{3} (2003), 3, 463--484.
\bibitem[15]{Haraux2009} R. Chill, A. Haraux, M. A. Jendoubi; \emph{Application of the \L{o}jasiewicz-Simon gradient inequality to gradient-like evolution equations}, Analysis and Applications, \textbf{7} (2009), 4, 351--372.
\bibitem[16]{Haraux2011} A. Haraux, M. A. Jendoubi; \emph{The \L{o}jasiewicz gradient inequality in the infinite-dimensional Hilbert space framework}, Journal of Functional Analysis, \textbf{260} (2011), 9, 2826--2842.
\bibitem[18]{Hussain_2005} J. Hussain; \emph{Analysis of some deterministic and stochastic evolution equations with solutions taking values in an infinite dimensional Hilbert manifold}, PhD thesis, University of York, 2015.
\bibitem[19]{Huss_2023} J. Hussain; \emph{Faedo-Galerkin approximations for nonlinear heat equation on Hilbert manifolds}, Carpathian Journal of Mathematics, \textbf{39}  (2023), 3, 667--682.
\bibitem[21]{Jendoubi} M. A. Jendoubi; \emph{A Simple Unified Approach to Some Convergence Theorems of L. Simon}, Journal of Functional Analysis, \textbf{153} (1998), 1, 187--202.
\bibitem[22]{Kapustyan} O. V. Kapustyan, D. V. Shkundin; \emph{Global Attractor of One Nonlinear Parabolic Equation}, Ukrainian Mathematical Journal, \textbf{55} (2003), 4, 535--547.
\bibitem[23]{Lions} P. L. Lions; \emph{Structure of the set of steady-state solutions and asymptotic behaviour of semilinear heat equations}, Journal of Differential Equations, \textbf{53}  (1984), 3, 362--386.
\bibitem[24]{Matano} H. Matano; \emph{Convergence of solutions of one-dimensional semilinear heat equations}, Journal of Mathematics of Kyoto University, \textbf{18} (1978), 2, 221--227.
\bibitem[26]{Rybka} P. Rybka; \emph{Convergence of Heat Flow On a Hilbert Manifold}, Proceedings of the Royal Society of Edinburgh: Section A Mathematics, \textbf{136}  (2006), 4, 851--862.
\bibitem[27]{Simon} L. Simon; \emph{Asymptotics for a Class of Non-Linear Evolution Equations, with Applications to Geometric Problems}, Annals of Mathematics, \textbf{118} (1983), 3, 525--571.
\bibitem[28]{Songu} Z. Songu; \emph{Non-Linear evolution equations}, Chapman \& Hall/CRC, 2004.
\bibitem[30]{Temam} R. Temam; \emph{Navier-Stokes Equations}, AMS Chelsea Publishing, 2000.
\bibitem[31]{Vishnevski} M. P. Vishnevskii; \emph{On stabilization of solutions to boundary value problems for quasilinear parabolic equations periodic in time}, Siberian Mathematical Journal, \textbf{34}  (1993), 5, 801--811.
\bibitem[32]{I.Vrabie} I. I. Vrabie; \emph{ $C_{0}$-Semigroups and Applications}, North-Holland Mathematics Studies, \textbf{191}, 2003.
\end{thebibliography}
\end{document}
